\documentclass{article}
\usepackage{amsmath}

\begin{document}

\section{Importancia de la medición y las medidas de tendencia central}

La medición es un proceso fundamental en la ciencia y en la vida cotidiana, ya que permite obtener información cuantitativa sobre diferentes fenómenos. Gracias a la medición, es posible comparar, analizar y tomar decisiones basadas en datos objetivos. Sin mediciones precisas, no sería posible el desarrollo de áreas como la física, la estadística o la ingeniería.

Dentro del análisis de datos, las medidas de tendencia central juegan un papel muy importante, ya que permiten resumir un conjunto de datos en un solo valor representativo. Estas medidas ayudan a comprender dónde se concentra la información y facilitan la interpretación de los resultados.

Una de las medidas más utilizadas es la media aritmética, que representa el promedio de un conjunto de datos. Se calcula sumando todos los valores y dividiendo entre el número total de observaciones:

$
\overline{x} = \frac{x_1 + x_2 + x_3 + \cdots + x_n}{n}
$

Otra medida importante es la mediana, que corresponde al valor central de un conjunto de datos ordenados. Si el número de datos es impar, la mediana es el valor que se encuentra en el centro; si es par, se calcula como el promedio de los dos valores centrales:

$
\text{Mediana} = \frac{x_{\frac{n}{2}} + x_{\frac{n}{2}+1}}{2}
$

En conclusión, las medidas de tendencia central son herramientas esenciales para resumir información y facilitar el análisis de datos, permitiendo una mejor comprensión de los fenómenos estudiados.

\end{document}